\subsection[Coarse K-homology and assembly]{Coarse $\boldsymbol{\K}$-homology and assembly}\label{sec:Assembly}

Using $\Gamma$-equivariant $\EE$-theory,\footnote{Of course, $\KK$-theory could also be used. In that case, a different description of the assembly maps would need to be used and hence completely new proofs of our results would be required. However, $\KK$-theory suffers from similar technical shortcomings as $\EE$-theory and therefore no advantage over our approach is expected. In the end, $\EE$-theory was primarily selected because of the author's greater familiarity with this theory.} we define the $\Gamma$-equivariant $\K$-homology with coefficients in a~separable $\Gamma$-\textCstar-algebra $D$ as the functors
\[
\K^\Gamma_*(\blank;D)\coloneqq \EE^\Gamma_*(\Cz(\blank);D).
\]
Note that this definition requires the function algebras $\Cz(\blank)$ to be separable, so we only consider these functors on the full subcategory of $\LCHp$ of all \emph{second countable} locally compact $\Gamma$-spaces.
In practice, this is only a minor restriction to the theory of Section~\ref{sec:Sigmastuff}, because we are mainly interested in two special cases: First, Rips complexes of discretizations of countably generated proper isocoarse $\Gamma$-spaces of bornologically bounded geometry are always second countable and, second, their compactification with a Higson dominated corona is second countable, if the corona is second countable, i.e., a compact metrizable space. The latter is an ubiquitous assumption in applications.


Due to the well-known properties of bivariant $\K$-theory (cf.\ \cite[Chapter~6]{GueHigTro}), the functors $\K^\Gamma_*(\blank;D)$ are $\Gamma$-equivariant single-space homology theories on the above-mentioned category. Furthermore, as direct limits preserve exact sequences, they can be extended to $\Gamma$-equivariant single-space homology theories on the full subcategory of $\sLCHp$ consisting of all second countable $\sigma$-locally compact $\Gamma$-spaces $\cX=\bigcup_{n\in\N}X_n$ by
\(\K^\Gamma_*(\cX;D)\coloneqq\lim_{n\to\infty}\K^\Gamma_*(X_n;D)\).


We want to compare the resulting transgression maps with the coarse assembly maps $\mu$: $\KX^\Gamma_*(X;D)\to\K_*(\Roe_\Gamma(X;D))$. The definition of the latter is based upon other pictures of $\K$-homology.
We choose the one using Yu's localization algebras \cite{YuLocalization} (see also \cite{QiaoRoe} and \cite{EngelWulffZeidler}) for our discussion.
A lot about localization algebras has been proven in the comprehensive book~\cite{WillettYuHigherIndexTheory}, but note that the ``localized Roe algebras'' considered there are slightly bigger than Yu's original localization algebras and hence a bit of care has to be taken when applying their results.

Following up our Section~\ref{sec:RoeAlg} about Roe algebras, let $X$ be a proper metric space equipped with a proper isometric $\Gamma$-action. We assume that sufficiently large $\Gamma$-$X$-module $(H,\rho,u)$ has been chosen and let $\HilbertMod\coloneqq H\otimes D$ be equipped with the induced representations. Recall that we defined $\Roe_\Gamma(X;D)$ as the norm closure in $\Lin(\HilbertMod)$ of the $*$-subalgebra of all $\Gamma$-equivariant locally compact operators of finite propagation, and if $A$ is a subspace of $X$, then $\Roe_\Gamma(A\subset X;D)$ denoted the ideal generated by all operators which are supported in an $R$-neighborhood of $A$ for some $R>0$ and $\Roe_\Gamma(X,A;D)\coloneqq \Roe_\Gamma(X;D)/\Roe_\Gamma(A\subset X;D)$.


\begin{defn}
%\label{defn:localizationalgebra}
The \emph{localization algebra} $\Loc(X;D)$ is the norm closure of the sub-$*$-algebra of $\Cb([1,\infty),\Roe_\Gamma(X;D))$ of all bounded uniformly continuous functions $L\colon[1,\infty)\to\Roe_\Gamma(X;D)$ such that the propagation of $L(t)$ is finite for all $t\geq 1$ and converges to zero as $t\to\infty$. Furthermore, we denote by $\Loc(A\subset X;D)\subset \Loc(X;D)$ the ideal generated by all $L$ for which there is a~function $\varepsilon\colon[1,\infty)\to(0,\infty)$ with $\varepsilon(t)\xrightarrow{t\to\infty}0$ such that $L(t)$ is supported in an $\varepsilon(t)$-neighborhood of $A$ for all $t\geq 1$, and we define $\Loc(X,A;D)\coloneqq \Loc(X;D)/\Loc(A\subset X;D)$.
\end{defn}

An obvious variant of \cite[Lemma~6.3.6]{WillettYuHigherIndexTheory}, which can be proven with the same methods, says that the inclusion $A\subset X$ induces via the theory of covering isometries an isomorphism $\K_*(\Loc(A;D))\cong\K_*(\Loc(A\subset X;D))$. Therefore we obtain a long exact sequence
\begin{align}
\cdots&\to\K_*(\Loc(A;D))\to \K_*(\Loc(X;D))\to \K_*(\Loc(X,A;D))\nonumber
\\&\to \K_{*-1}(\Loc(A;D))\to\cdots.\label{eq:localizationLES}
\end{align}


The relation between localization algebras and our definition of $\K$-homology is the following. As $\Kom(\HilbertMod)\cong D\otimes\Kom(H)$, it is straightforward to verify (cf.\ \cite[Corollary~4.2]{QiaoRoe}) that
\begin{align*}
\delta\colon\ \Loc(X;D)\maxtensor\Cz(X) &\to\frac{\Cb([1,\infty),D\otimes \Kom(H))}{\Cz([1,\infty),D\otimes\Kom(H))},
\\
L\otimes f&\mapsto [t\mapsto L(t)\circ (\rho(f)\otimes\id_D)]
\end{align*}
defines an $\Gamma$-equivariant asymptotic morphism. It vanishes on the ideal $\Loc(A\subset X;D)\maxtensor\Cz(X\setminus A)$ and hence we also obtain an induced $\Gamma$-equivariant asymptotic morphism
\[
\delta\colon\ \Loc(X,A;D)\maxtensor\Cz(X\setminus A) \to\frac{\Cb([1,\infty),D\otimes \Kom(H))}{\Cz([1,\infty),D\otimes\Kom(H))}.
\]
Both of them define canonical $\EE^\Gamma$-theory elements in $\EE^\Gamma_0(\Loc(X;D)\maxtensor\Cz(X),D)$ and \linebreak $\EE^\Gamma_0(\Loc(X,A;D)\maxtensor\Cz(X\setminus A),D)$, respectively, which we shall also denote by the letter~$\delta$.\footnote{Here we need $\EE^\Gamma$-theory also for non-separable \textCstar-algebras. It can be defined in an ad hoc way as in \cite[Section~1]{WulffTwisted}, or even simpler by letting $\EE^\Gamma_*(A,B)$ for possibly non-separable $A$, $B$ be obtained from the groups $\EE^\Gamma_*(A',B')$ by taking the direct limit over all separable $B'\subset B$ and the inverse limit over all separable $A'\subset A$. The relevant maps considered in the remainder of this section will be the same for both versions.
Note that these generalizations retain the composition and tensor products of $\EE^\Gamma_*$-theory, but be aware that they are simply too naive to preserve the long exact sequences. }
We obtain the composed group homomorphism
\begin{align*}
\Delta\colon\ \K_*(\Loc(X;D))&\cong\EE_*(\C,\Loc(X;D))
\\&\xrightarrow{\blank\otimes\id_{\Cz(X)}}\EE^\Gamma_*(\Cz(X),\Loc(X;D)\maxtensor\Cz(X))
\\&\xrightarrow{\delta\circ\blank}\EE^\Gamma_*(\Cz(X),D)=\K^\Gamma_*(X;D)
\end{align*}
and similarily $\Delta\colon \K_*(\Loc(X,A;D))\to\K^\Gamma_*(X,A;D)$ in the relative case.
By exploiting the homological properties of domains and targets of the maps $\Delta$ (in particular that the domains are functorial under uniformly continuous coarse $\Gamma$-maps via the theory of covering isometries, that the maps $\Delta$ are natural under such maps and that the connecting homomorphisms are in both cases constructed via mapping cones) we see that the maps $\Delta$ map \eqref{eq:localizationLES} to the long exact sequence of $\K^\Gamma_*(\blank,D)$.

\begin{thm}\label{thm:Khomologydefinitions}
The maps $\Delta$ are isomorphisms in the absolute case for every proper metric space~$X$ equipped with a proper isometric $\Gamma$-action and hence by a five lemma argument also in the relative case.\footnote{Instead of Theorem~\ref{thm:Khomologydefinitions}, preprint versions of this article contained the ``Assumption 6.8'' that $\Gamma$ and $D$ have been chosen for which the statement of this theorem holds. Thanks to a referee's hint, we can now prove it in full generality.}
\end{thm}

Proofs of some special cases of the theorem have been known for a long time.
In the non-equivariant case ($\Gamma=1$) without coefficients ($D=\C$), one proof is sketched in \cite[Construction~6.7.5, Exercises~B.3.2 and~B.3.3]{WillettYuHigherIndexTheory}, albeit for a different version of the localization algebra. Note, however, that this proof relies solely on homological properties and a comparison on one-point spaces. Therefore, it goes through for original version of the localization algebras, too, and even with arbitrary coefficient \textCstar-algebras $D$.

Another proof for the case $\Gamma=1$ and $D=\C$ can be found in \cite[Proposition~4.3]{QiaoRoe} and relies on Paschke duality. The author is indebted to a referee for pointing out a version of Paschke duality for the equivariant $\KK$-theory groups $\KK^\Gamma_*(\Cz(X),D)$ which allows to prove the theorem in full generality.
Nevertheless, as localization algebras and $\EE$-theory are defined in a very similiar manner, it would be interesting to know if the theorem can be proven directly without taking the detour via Paschke duality.

\begin{proof}
It suffices to prove the theorem for separable $D$, because the general case follows by taking direct limits over all separable sub-\textCstar-algebras.
In this case, we can make use of the Paschke duality for equivariant $\KK$-theory that was discussed in \cite[Appendix~B]{GueWillYu_DyncomplKtheory}.

An adjointable operator $T$ on $\HilbertMod$ is called pseudolocal if the commutators $\rho(f)T-T\rho(f)$ are compact for all $f\in \Cz(X)$.
Let $\PseudoLoc_\Gamma(X;D)$ denote the sub-\textCstar-algebra of $\Lin(\HilbertMod)$ generated by the $\Gamma$-equivariant pseudolocal operators that have finite propagation and let $\PseudoLocLoc(X;D)$ denote the sub-\textCstar-algebra of $\Cb([1,\infty),\PseudoLoc_\Gamma(X;D))$ of all bounded uniformly continuous functions $L\colon[1,\infty)\to\PseudoLoc(X;D)$ such that the propagation $L(t)$ is finite for all $t\geq 1$ and converges to zero as $t\to\infty$.\looseness=-1

We are free to choose a very specific sufficiently large $\HilbertMod=H\otimes D$, namely $\HilbertMod=\HilbertMod_D\coloneqq H'\otimes \ell^2(\N)\otimes\ell^2(\Gamma)\otimes D$ for a non-degenerate $\Gamma$-$X$-module $(H',\rho',u')$.
It will be equipped with $\Gamma$-action $\varepsilon$ which acts diagonally on the first, third and fourth factor and the representation $\pi$ of $\Cz(X)$ which acts on the first factor.
For this choice, \cite[Appendix B]{GueWillYu_DyncomplKtheory} contains the proof of isomorphisms
\begin{gather*}
\KK_p^\Gamma(\Cz(X),D)\cong\K_{p+1}\biggl(\frac{\PseudoLoc_\Gamma(X;D)}{\Roe_\Gamma(X;D)}\biggr) \cong\K_{p+1}\biggl(\frac{\PseudoLocLoc(X;D)}{\Loc(X;D)}\biggr)\cong\K_p(\Loc(X;D))
\end{gather*}
for separable $\Gamma$-\textCstar-algebras $D$ and proper metric spaces $X$ equipped with a proper cocompact isometric $\Gamma$-action.
